\scanpage{12}
2. \uline{Fonctions de Weyl}

\margin{\(K^\infty,\ L^\infty(K)\)\par \unsure{fct} \unsure{définies} \unsure{sur} \(K^\infty_K\)}

\uline{Définition 1} Soit \deleted{\(M\) \unsure{un espace mesuré} \unsure{(mesure} \(m)\)} \inserted{\(I=]0,\omega[\quad(0<\omega\le+\infty)\)}. On appelle \uline{fonction de Polya} sur \deleted{\(\unsure{M}\)} une fonction \(P\) réelle définie sur l'ens \(K^\infty(\deleted{\unsure{M}}\inserted{I})\) des fonctions mesurables à support compact, ayant les propriétés suivantes :

\figure{\begin{tikzpicture}\draw (0,0)--(0,4.8);\end{tikzpicture}}

\(P_1)\quad P\) est \uline{symétrique}, i.e. \(P(f)=P(g)\) si \(f\) et \(g\) sont spectralement équivalentes.

\(P_2)\quad P\) est convexe, i.e.
\[
P(\lambda f+(1-\lambda)g)\le\lambda P(f)+(1-\lambda)P(g)
\]
si \(0\le\lambda\le1\).

\(P_3)\quad\) Sur \deleted{tout} compact \(K\), \deleted{\unsure{P est}} et toute suite \((f_i)\) dans \(\unsure{K^\infty_K(K)|K^\infty_K(I)}\) tendant vers une \deleted{\unsure{limite}} limite \(f\) en mesure, on a
\[
P(f)\le\underline{\lim}\,P(f_i).
\]

En fait, \(P_3\) est équivalente à la condition plus forte en apparence

\(P'_3)\quad\) \inserted{Pour tout compact \(K\)} \(P\) est semi-continue inférieurement sur \deleted{tout \(K\)} pour la topologie faible sur \(L^\infty(K)\).

Il faut montrer que pour tout réel \(a\), \deleted{\(\unsure{P^{-1}(]-\infty,a[)}\)} \inserted{l'ens \(A\) des} \deleted{est \unsure{fblt} fermé} \(f\in L^\infty(K)\) telles que \(P(f)\le a\) est \unsure{fblt} fermé.

\deleted{\unsure{D'après} \illegible{2} \unsure{de Banach-Dieudonné} \(\unsure{A}\) \illegible{2}}

\deleted{\unsure{Il suffit de prouver} \unsure{que sur} \unsure{toute partie} \illegible{1} \unsure{de} \unsure{tout} \unsure{ensemble} \unsure{fermé} \unsure{de} \(L^\infty(K)\) \unsure{est fblt} \illegible{2}}

\inserted{Cet ensemble est convexe en} vertu de \(P_2\), \deleted{\unsure{et par conséquent}} \inserted{et en vertu de \(P_3\) il est} fermé pour \unsure{une suite bornée|des suites bornées} dans \(L^\infty(K)\) qui \unsure{converge|convergent} en mesure. Or un tel ensemble \(A\) est déjà \unsure{fblt} fermé, car en vertu du th. de Banach-Dieudonné, on est ramené au cas où \(A\) est borné dans \(L^\infty(K)\). Il est donc \unsure{fblt relt cpt} dans \(L^{\unsure{1|2}}(K)\), \deleted{\unsure{donc}} de plus il est \uline{fermé} dans \(L^1(\unsure{K|k})\) (pour la top. forte, donc pour la

\mspage{13}
